% GENERATED FILE. Edit canonical lessons, then run npm run generate:books.
% canonical-source-hash: fnv1a64:bb29e84a4dd2fe79
% canonical-lessons: TE-C85-cold, TE-C85-heat, TE-C85-fear, TE-C85-anger, TE-C85-sorrow

\chapter{Cold, Hot, and How You Feel}
\label{ch:te-cold-hot-feel}
\hlchaptermodality{85}

\begin{chapteropening}
\textbf{By the end of this chapter:} \emph{I can say that I feel cold or hot, and that I am afraid, angry or sad.}

Complete TE-C85-sorrow: say how you feel, from the weather to the heart.
\end{chapteropening}

\section[cali]{\te{చలి} (cali) — cold, the cold}

\label{lesson:TE-C85-cold}



\begin{quote}
Before the new one: say the Telugu for thirst, then the Telugu for rest.
\end{quote}

\subsection*{You'll want to know: \te{చలి}}
\textbf{\te{చలి}} (\emph{cali}) — \textquotedblleft{}cold, the cold\textquotedblright{}.

\textbf{\te{ఇవాళ} \te{చలిగా} \te{ఉంది}} — \textquotedblleft{}it is cold today\textquotedblright{}. \textbf{\te{నాకు} \te{చలిగా} \te{ఉంది}} — \textquotedblleft{}I feel cold\textquotedblright{}.

The same frame again: with \te{నాకు} it is your body, without it it is the day.

\begin{grammarlens}[title={What you've built}]
The first word that works for both the weather and the body.
\end{grammarlens}

\subsection*{Your turn}
\begin{itemize}
\raggedright
  \item \emph{Say it:} \emph{cali}
  \item \emph{Say it:} \emph{cali}, once more
  \item \emph{Say it:} say \emph{ivāḷa caligā undi}, then \emph{nāku caligā undi}
  \item \emph{Recall:} say the Telugu for thirst, then the Telugu for rest, then say \emph{cali} again
\end{itemize}

\begin{tcolorbox}[breakable,colback=teal!4,colframe=teal!35!black,title={Before you move on}]
What does \te{చలి} mean? (\textquotedblleft{}Cold, the cold\textquotedblright{}.) Say it once more.
\end{tcolorbox}
\section[vēḍi]{\te{వేడి} (vēḍi) — heat; hot}

\label{lesson:TE-C85-heat}



\begin{quote}
Before the new one: say the Telugu for rest, then the Telugu for the cold.
\end{quote}

\subsection*{You'll want to know: \te{వేడి}}
\textbf{\te{వేడి}} (\emph{vēḍi}) — \textquotedblleft{}heat\textquotedblright{}, and before a noun, \textquotedblleft{}hot\textquotedblright{}.

\textbf{\te{ఇవాళ} \te{చాలా} \te{వేడిగా} \te{ఉంది}} — \textquotedblleft{}it is very hot today\textquotedblright{}. \textbf{\te{వేడి} \te{నీళ్ళు}} — \textquotedblleft{}hot water\textquotedblright{}.

It is the pair of \te{చలి}, and it slots into the same frame.

\begin{grammarlens}[title={What you've built}]
Cold and hot, as a pair, in one frame.
\end{grammarlens}

\subsection*{Your turn}
\begin{itemize}
\raggedright
  \item \emph{Say it:} \emph{vēḍi}
  \item \emph{Say it:} \emph{vēḍi}, once more
  \item \emph{Say it:} say \emph{vēḍi nīḷḷu}, then \emph{ivāḷa cālā vēḍigā undi}
  \item \emph{Recall:} say the Telugu for rest, then the Telugu for the cold, then say \emph{vēḍi} again
\end{itemize}

\begin{tcolorbox}[breakable,colback=teal!4,colframe=teal!35!black,title={Before you move on}]
What does \te{వేడి} mean? (\textquotedblleft{}Heat; hot\textquotedblright{}.) Say it once more.
\end{tcolorbox}
\section[bhayaṁ]{\te{భయం} (bhayaṁ) — fear}

\label{lesson:TE-C85-fear}



\begin{quote}
Before the new one: say the Telugu for the cold, then the Telugu for heat.
\end{quote}

\subsection*{You'll want to know: \te{భయం}}
\textbf{\te{భయం}} (\emph{bhayaṁ}) — \textquotedblleft{}fear\textquotedblright{}.

\textbf{\te{నాకు} \te{భయంగా} \te{ఉంది}} — \textquotedblleft{}I am afraid\textquotedblright{}. The frame that carried tiredness and cold now carries a feeling, with nothing else changed.

\begin{grammarlens}[title={What you've built}]
The first feeling, in a frame you can already build.
\end{grammarlens}

\subsection*{Your turn}
\begin{itemize}
\raggedright
  \item \emph{Say it:} \emph{bhayaṁ}
  \item \emph{Say it:} \emph{bhayaṁ}, once more
  \item \emph{Say it:} say \emph{nāku bhayaṁgā undi}
  \item \emph{Recall:} say the Telugu for the cold, then the Telugu for heat, then say \emph{bhayaṁ} again
\end{itemize}

\begin{tcolorbox}[breakable,colback=teal!4,colframe=teal!35!black,title={Before you move on}]
What does \te{భయం} mean? (\textquotedblleft{}Fear\textquotedblright{}.) Say it once more.
\end{tcolorbox}
\section[kōpaṁ]{\te{కోపం} (kōpaṁ) — anger}

\label{lesson:TE-C85-anger}



\begin{quote}
Before the new one: say the Telugu for heat, then the Telugu for fear.
\end{quote}

\subsection*{You'll want to know: \te{కోపం}}
\textbf{\te{కోపం}} (\emph{kōpaṁ}) — \textquotedblleft{}anger\textquotedblright{}.

\textbf{\te{నాకు} \te{కోపం} \te{వచ్చింది}} — \textquotedblleft{}I got angry\textquotedblright{}, literally \textquotedblleft{}to me anger came\textquotedblright{}. Anger arrives the way a fever did, on the come-verb, and the person it happens to sits in the dative.

\begin{grammarlens}[title={What you've built}]
A feeling that comes to you, like a fever.
\end{grammarlens}

\subsection*{Your turn}
\begin{itemize}
\raggedright
  \item \emph{Say it:} \emph{kōpaṁ}
  \item \emph{Say it:} \emph{kōpaṁ}, once more
  \item \emph{Say it:} say \emph{nāku kōpaṁ vaccindi}
  \item \emph{Recall:} say the Telugu for heat, then the Telugu for fear, then say \emph{kōpaṁ} again
\end{itemize}

\begin{tcolorbox}[breakable,colback=teal!4,colframe=teal!35!black,title={Before you move on}]
What does \te{కోపం} mean? (\textquotedblleft{}Anger\textquotedblright{}.) Say it once more.
\end{tcolorbox}
\section[bādha]{\te{బాధ} (bādha) — sorrow, distress}

\label{lesson:TE-C85-sorrow}



\begin{quote}
Before the new one: say the Telugu for fear, then the Telugu for anger.
\end{quote}

\subsection*{You'll want to know: \te{బాధ}}
\textbf{\te{బాధ}} (\emph{bādha}) — \textquotedblleft{}sorrow, distress\textquotedblright{}.

\textbf{\te{నాకు} \te{బాధగా} \te{ఉంది}} — \textquotedblleft{}I feel sad\textquotedblright{}. It also covers the ache behind worry, so it answers \emph{\te{మీరు} \te{ఎలా} \te{ఉన్నారు}?} honestly on a bad day.

\begin{grammarlens}[title={What you've built}]
An honest answer to how are you, when the answer is not \te{బాగా}.
\end{grammarlens}

\subsection*{Your turn}
\begin{itemize}
\raggedright
  \item \emph{Say it:} \emph{bādha}
  \item \emph{Say it:} \emph{bādha}, once more
  \item \emph{Say it:} say \emph{nāku bādhagā undi}
  \item \emph{Recall:} say the Telugu for fear, then the Telugu for anger, then say \emph{bādha} again
\end{itemize}

\begin{tcolorbox}[breakable,colback=teal!4,colframe=teal!35!black,title={Before you move on}]
What does \te{బాధ} mean? (\textquotedblleft{}Sorrow, distress\textquotedblright{}.) Say it once more.
\end{tcolorbox}
